%2. Resource changes
%Figure 4 is most informative in the comparison between the bottom 40 percent of treated municipalities and control municipalities, which experienced similar resource growth. The interpretation should nevertheless be cautious. In particular, not even the bottom group or the control group experience zero resource growth. The exercise therefore does not isolate the effects of the geographic reorganization, absent those of changes in resources. Rather, it shows that the estimated gains are not confined to municipalities experiencing the largest resource expansions.
%
%Relatedly, I think Table B3 is helpful. I would not dismiss these specifications simply because resource changes are themselves induced by treatment. Conditioning on policy-induced resource changes can still be informative about whether the reduced-form effect is concentrated in places experiencing large resource expansions.


%Internal note: No he cambiado el paper para responder a esta pregunta. Seria quitar la footnote y poner en su lugar el texto que se pone mas abajo. El apendice es igual.

Thank you very much for these comments.

On the comparison between municipalities with control-like versus larger resource growth: we agree with your characterization and do not read this exercise as isolating the effect of the geographic reorganization from that of resource changes. Our point is narrower. Finding sizeable effects even among municipalities whose resource growth resembles that of the control group does not mean that resource growth was unimportant---resources may well have been necessary to enable the effective implementation of the deployment reorganization even in that municipalities---but it does indicate that the reorganization itself plays an important role. This qualification is in the updated paper. For example, in the Introduction, we state that ``these patterns indicate that while resource increases helped enable implementation and may explain part of the effects in some municipalities, the beat-policing components of the reform were central drivers of the overall gains.'' In Section \ref{mechanisms}, we similarly state that ``although we cannot fully separate the beat policing component from the resource component of the reform---resources may have been important to enable effective implementation of the program, or may explain part of the effects in some municipalities---the fact that municipalities with resource increases similar to those experienced by control municipalities still display sizeable effects indicates that the effects of \emph{Plan Cuadrante} are not driven entirely by resources.''

Regarding Table B3: thank you for the suggestion; we agree these estimates are informative. This is now discussed directly in the main text of Section \ref{mechanisms}, rather than only in a footnote, in addition to the Online Appendix. 

We included the following text in section 4.4: ``As a further check, we re-estimate the results in Figure 2 controlling directly for police resources. Online Appendix Table \ref{tab:results_controlbyresources} shows that the magnitude of the estimated effects remains similar for most variables, although the statistical significance is reduced in some estimations. This likely reflects a mechanical feature of the \citet{Callaway} estimation procedure: control variables are not incorporated parametrically, but are instead used to construct the comparison group via outcome regression or propensity-score weighting, so the resulting estimation uncertainty propagates to the ATT standard errors. Note, however, that police resources are probably a bad control---i.e., at least part of their variation is a result of the adoption of \emph{Plan Cuadrante}---so these results should be interpreted cautiously.'' 


The Online Appendix (Section \ref{resources_plan_cuadrante}) reads: ``The results of this analysis are reported in Table \ref{tab:results_controlbyresources}. The magnitude of the estimated effects remains similar for most variables, although the statistical significance is reduced in some estimations. This likely reflects a mechanical feature of the \citet{Callaway} estimation procedure: control variables are not incorporated parametrically, but are instead used to construct the comparison group via outcome regression or propensity-score weighting, so the resulting estimation uncertainty propagates to the ATT standard errors. Note, however, that these resource variables are bad controls, as they are affected by the adoption of \emph{Plan Cuadrante}, so these results should be interpreted with caution.''

